There are over 13,000 school districts that operate in the United States run independently by school boards. They are responsible for annual expenditures  the size of the US Department of Defense budget. These decisions ripple through every American community annually and have wide reaching implications for everything from property values to economic development to the content and quality of education provided to students in each community. Despite this wide reaching scope and the importance of these decisions, local school boards remain woefully understudied (Land 2002). Limited work has been done on exploring the conditions under which voters participate in school board elections and candidates choose to run, but never with an eye on evaluating causal relationships.

This dissertation focuses on a single critical piece of the local election puzzle--how does an extremely sudden, contentious, and deeply partisan state level policy shock affect local governments? Specifically:
\vspace{-6mm}
\begin{enumerate}
\begin{singlespace}
\setlength\itemsep{-6pt}
\item How does such a policy shock affect the level of participation in school board elections for voters and challengers to incumbent board members?
\item How does that shock change school board policymaking?
\item Does sudden polarization of state level politics lead to evidence of partisanship on school boards?
\end{singlespace}
\end{enumerate}
\vspace{-1mm}

The recent reforms enacted by Governor Scott Walker in Wisconsin have created a natural experiment within the state that allows causal leverage on these questions for the first time. Boards were given unanticipated freedom in setting employee compensation and work rules, and education became the premier issue in a bitterly contested statewide series of elections. By comparing the behavior of candidates, board members, boards, and voters before and after the policy shock it is possible to evaluate the causal impact of state level policy on the politics of school boards. Finally, by tying board policy, electoral activity, and public policy preferences together along a single issue dimension identified exogenously to all school boards a spatial model of school board elections can be evaluated empirically for the first time.

Using results from the universe of school district elections over the last 6 years, and combining them with results from elections for statewide office, I can evaluate the impact of state level exogenous policy shock on voter and candidate participation in school board elections. Additionally, by collecting data on school board policy choices in the newly unconstrained policy space of employee compensation, the issue preferences of school boards can be examined. 